\documentclass[11pt,a4paper]{article}
\usepackage[utf8]{inputenc}
\usepackage[T1]{fontenc}
\usepackage[romanian]{babel}
\usepackage[margin=2.2cm]{geometry}
\usepackage{amsmath,amssymb}
\usepackage{enumitem}
\usepackage{parskip}
\pagestyle{empty}

\begin{document}
\begin{center}
{\large\bfseries MODEL DE ANTRENAMENT --- TS2}\\[2pt]
{\small (subiect generat pentru autoevaluare, nu e biletul oficial de examen)}
\end{center}
\vspace{0.3cm}

\begin{enumerate}[leftmargin=1.4em]

\item Se consideră sistemul reprezentat prin următoare funcție de transfer:
$$G_0(s)=\frac{2s+3}{s^2+7s+10}$$
\begin{enumerate}[label=\alph*)]
\item Să se determine modelul la stare. (0.75p)
\item Să se realizeze schema bloc. (0.5p)
\item Să se determine răspunsul procesului (la intrare treaptă unitară). (0.75p)
\item Să se determine dacă procesul este controlabil. (0.5p)
\end{enumerate}

\item Fie un sistem cu reacție negativă unitară în care:
$$G(s)=8\,\frac{0.25s+1}{s}$$
\begin{enumerate}[label=\alph*)]
\item Trasați caracteristicile amplitudine-pulsație și fază-pulsație. (1p)
\item Să se calculeze eroarea staționară. (0.5p)
\end{enumerate}

\item Fie un sistem de reglare cu funcția de transfer pe calea directă $G(s)$ și funcția de transfer pe calea de reacție, $H(s)$:
$$G(s)=\frac{k}{s(s+4)}\qquad H(s)=\frac1{s+2}$$
\begin{enumerate}[label=\alph*)]
\item Folosind curba polară să se determine domeniul de valori al factorului de amplificare pentru care sistemul în circuit închis este stabil. Trasați curba polară. (1p)
\item Să se calculeze marginea de fază. (0.5p)
\end{enumerate}

\item Se dă sistemul cu reacție negativă unitară, cu funcția de transfer în circuit deschis:
$$G(s)=\frac{K(s+4)}{s(s+3)(s^2+2s+5)}$$
Să se traseze locul rădăcinilor. (1.5p)

\item Fie sistemul în circuit închis de următoarea formă:
$$G(s)=\frac{16}{s^2+8s+16}$$
Să se calculeze și să se schițeze răspunsul sistemului pentru o mărime de intrare $u(t)=3\sin3t$. (1p)

\end{enumerate}

\textbf{Teorie:}
\begin{enumerate}[leftmargin=1.4em]
\item Cum se aproximează, în domeniul frecvenței, caracteristica amplitudine-pulsație pentru un element de ordinul doi cu poli complex conjugați?
\item Ce reprezintă frecvența de frângere $\omega_f$ a unui factor de ordinul întâi și cum modifică ea panta asimptotei Bode?
\item Care sunt indicatorii de calitate din domeniul timpului ce caracterizează un sistem în circuit închis?
\item Să se realizeze schema bloc de reglare pentru un sistem cu reacție negativă unitară și element de execuție cu funcție de transfer $H_e(s)$.
\end{enumerate}

\vspace{0.2cm}
\small Toate subiectele teoretice valorează 1 punct.

\end{document}
